% ============================================================
% Page 4 — Alban du IVème siècle à Saint Benoît
% Transcription validée du document source.
% Mise en page simple ; typographie et lettrines à traiter
% ultérieurement dans le préambule et le modèle général.
% ============================================================

\thispagestyle{plain}

\vspace*{12mm}

\begin{center}
    {\Large\bfseries Alban du IVème siècle à Saint Benoît}
\end{center}

\vspace{18mm}

\noindent
Il est normal de penser qu'Alban fut christianisé d'assez bonne heure.
Le christianisme prenait naissance en ville avant de s'étendre à la
campagne. Le mot latin « paganus » a donné en français le doublet
« paysan » et « païen ». L'Albigeois ayant été évangélisé au IVème
siècle par Saint Clair, Saint Amarand, on peut penser que le relais
gallo-romain d'Alban a formé un petit noyau chrétien, peut-être autour
de cet « ALBANUS » qui aurait accolé son nom à celui de Notre Dame.
Les spécialistes de l'histoire religieuse prétendent que les paroisses
sont les plus anciennes -- qui se sont mises sous le vocable de Notre
Dame. C'est le cas d'Alban qui deviendra « Sainte Marie d'Alban ».

\par\vspace{12mm}

\noindent
Les invasions barbares affectent-elles notre petite agglomération ?
La réponse se perd dans les brumes de ces lointaines époques. Un
couvercle de sarcophage dit « Wisigoth » (actuellement appuyé à la
façade ouest de l'église) semble indiquer que le contrecoup de nouveaux
arrivants affecta nos régions. L'intrusion des Arabes dans le midi de
la France y a laissé des traces avec la culture du sarrasin ou blé noir.
Et çà et là quelques individus rappellent par leur teint particulièrement
bronzé et des cheveux très noirs -- le type arabe -- Il est difficile,
à l'heure présente, d'en dire plus.

\vfill
